\chapter{Variational Calculus}
\label{chap:variational-calculus}

% Closed question: What selects one admissible continuation as the history
% actually carried forward?

\begin{chapteressay}

A bead slides down a wire from a higher point $A$ to a lower point
$B$. The wire's shape is to be designed for fastest descent. The
straight line connecting $A$ and $B$ is not the answer; a curve that
drops more steeply at first builds up speed faster. The optimal
curve is the cycloid: the path traced by a point on the rim of a
rolling wheel.

The cycloid is not obvious. It is also not arbitrary. It is forced
by a precise selection principle: among all admissible curves
connecting $A$ to $B$, the cycloid is the unique one minimizing the
travel time. The selection is unique; the cycloid is the answer.

This is the chapter's structural question. Given data (the endpoints
$A$, $B$) and an admissible family of continuations (smooth curves
between them), how is one continuation selected as the actual
history? The answer is variational: select the continuation that
extremizes a chosen functional. For the brachistochrone, the
functional is the descent time; the extremum is a minimum; the
selected curve is the cycloid.

The chapter develops the selection principle through five
structures. The first is minimality as selection: among
admissible alternatives, the minimum of the functional is the
selected element. The soap film example illustrates this: the
film selects the minimum-area surface for the given boundary.

The second is the Euler-Lagrange equation. Given a functional
$J[y] = \int L(y, y', x) \, dx$, the Euler-Lagrange equation
$\partial L/\partial y - d/dx \, \partial L/\partial y' = 0$ is
the necessary condition for $y$ to be an extremum. The
brachistochrone problem is solved by this technique; the cycloid
is the unique smooth curve satisfying the boundary conditions
and the Euler-Lagrange equation. This is Medium-section
material: the derivation requires careful variational
calculation.

The third is Kolmogorov selection. For discrete settings, the
selection principle ``minimize $K$'' is conceptually parallel
to the variational principle but algorithmically different.
Exact $K$ is not computable; the discipline uses computable
proxies (heartbeat counts, BIC, integrated curvature, edit
distance). This is also Medium-section material: the connection
between continuous variational calculus and discrete
algorithmic minimization is the load-bearing conceptual
content.

The fourth is fluxions resolved. Newton's infinitesimal calculus
is reframed as the limit of finite-difference ratios. The
derivative is a Galerkin limit, not a ghost of departed
quantities. The Euler-Lagrange equation operates on these
limits without requiring the infinitesimal vocabulary.

The fifth is Galerkin convergence. The discrete Galerkin
approximations to a variational problem converge to the
continuous extremum as the discretization refines. The
convergence unifies the discrete and continuous: both are
expressions of the same selection principle.

The chapter's ground example is the brachistochrone. The cycloid
is the unique minimum-time path from $A$ to $B$ under gravity.
The Euler-Lagrange equation produces the cycloid; the cycloid is
the selected history.

The chapter's second concrete example is the calligrapher
drawing one stroke. The calligrapher's brush travels from a
specific starting point to a specific ending point. Among all
the trajectories the brush could take, the calligrapher's
training selects the unique minimum-effort trajectory: the
trajectory that produces the desired stroke without
unnecessary inflection or wasted motion. The selection is the
calligrapher's body's implementation of the chapter's
variational principle. After thirty years of practice, the
selection is automatic; but it is the chapter's principle being
performed.

The chapter's closed question is:

\begin{center}
\emph{What selects one admissible continuation as the history
actually carried forward?}
\end{center}

The chapter's answer is: a variational principle. Among
admissible continuations, the one that extremizes a chosen
functional is the selected history. The functional may be the
energy, the action, the surface area, the descent time, the
Kolmogorov complexity, or any other admissible functional. The
Euler-Lagrange equation is the necessary condition for the
extremum. The discrete version uses computable proxies for the
incomputable exact selection.

The chapter is the central analytic chapter of the book, paired
with Chapter 5. Chapter 5 established the discrete (spline)
case; Chapter 6 extends to continuous (variational) selection.
Together they give the mathematical gauge its selection tools.

Subsequent chapters build on this. Chapter 7 introduces
transport: how information moves between selected histories.
Chapter 8 introduces calibration: how the reference structure is
made commensurable across observers. Chapter 9 introduces
curvature: the geometric quantity recording the failure of
transport to commute. Chapter 10 introduces invariance: what
survives admissible relabeling. Chapter 11 closes the book with
measure and entropy.

\end{chapteressay}

\section{Minimality as Selection Principle}
\label{se:minimality-as-selection-principle}

A soap film stretched across a wire frame takes a specific shape: the
shape that minimizes the surface area subject to the boundary
condition of the frame. The film does not assume an arbitrary shape;
it does not undulate; it does not choose any of the infinitely many
surfaces that fit the wire boundary. The film selects the unique
minimal-area surface.

The shape is forced by physics. The film's surface tension is a force
proportional to area. The film's energy is the surface tension times
the area. The film equilibrates at the minimum-energy configuration:
the surface of minimum area for the given boundary. The selection is
a consequence of the energy minimization.

Mathematically, the minimal surface satisfies the minimal surface
equation:
\[
\nabla \cdot \left(\frac{\nabla u}{\sqrt{1 + |\nabla u|^2}}\right) = 0,
\]
where $u(x, y)$ is the surface height at horizontal position $(x, y)$.
The equation is non-linear in general, but for small slopes it
reduces to the linear Laplace equation $\Delta u = 0$. The surface
is harmonic; among all harmonic functions matching the boundary
values, it is unique.

This is the chapter's first claim. Minimality is a selection principle.
Among all admissible continuations consistent with the boundary data,
the minimum (of some chosen energy functional) is uniquely selected.
The selection is forced by the functional, not chosen arbitrarily.

The previous chapter introduced the spline as the minimum-curvature
interpolant of finite data. The current chapter extends this to
selection problems: given an admissible family of continuations
parameterized by initial conditions, which member of the family is
the actual history? The answer is the member that minimizes some
functional.

The functional is the chapter's algebraic object. Common functionals
include:
\begin{itemize}
\item Surface area for soap films.
\item Travel time for light rays in optical media.
\item Action (kinetic minus potential energy integrated over time)
for mechanical paths.
\item Information content (Kolmogorov complexity) for sequences.
\item Integrated curvature for splines.
\end{itemize}
Each functional selects a unique admissible continuation. The
specific functional depends on the physical situation; the selection
principle (minimize the functional) is universal.

In Lean, the selection principle appears throughout. The
\texttt{GalerkinProcess} typeclass from Chapter 5 selects the
best approximation in the energy norm. The \texttt{Spline.le}
relation orders splines by their refinement: among admissible
splines, the most refined consistent with the data is the
minimum-curvature interpolant. The compiler's elaborator selects
the minimum-effort proof when multiple proofs are admissible.

The chapter's discipline: among admissible continuations, the
minimum of the chosen functional is the selected one. The
selection is unique when the functional is strictly convex. The
selection may be non-unique for non-convex functionals; in that
case, additional criteria are needed to break the tie.

The minimal surface example illustrates the existence and uniqueness
theory. For a smooth wire frame (the boundary), the minimal surface
exists (by the direct method of the calculus of variations) and is
unique (when the boundary is sufficiently regular). The proof
relies on the convexity of the surface area functional.

For non-convex functionals --- the chapter's example will be the
action functional for mechanical paths --- existence and uniqueness
are more delicate. The Euler-Lagrange equation (next section) gives
the necessary condition for an extremum; whether the extremum is a
minimum, maximum, or saddle point depends on second-order
information.

\section{The Euler-Lagrange Equation}
\label{se:the-euler-lagrange-equation}

The brachistochrone problem, posed by Johann Bernoulli in 1696, asks:
given two points $A$ (higher) and $B$ (lower, not directly below $A$),
which curve connecting them produces the fastest descent of a bead
sliding under gravity, starting from rest at $A$?

The straight line is not the answer. The straight line minimizes
distance; it does not minimize time. A curve that drops more steeply
at first builds up speed faster; the higher speed compensates for
the longer path. The optimal curve is therefore neither the
straightest nor the shortest; it is the curve that balances distance
and acceleration to minimize total time.

The answer is the cycloid: the curve traced by a fixed point on the
rim of a wheel rolling along a horizontal line. Parameterized by the
wheel's angle $\theta$:
\[
x(\theta) = r(\theta - \sin\theta), \quad y(\theta) = -r(1 - \cos\theta),
\]
where $r$ is the wheel's radius. The cycloid is the unique
admissible curve through $A$ and $B$ that minimizes the time of
descent.

The proof is a paradigmatic application of the calculus of variations.

\emph{Setup.} Let $y(x)$ be the curve from $A = (0, 0)$ to $B = (a,
-b)$ with $b > 0$. By energy conservation, a bead at height $-y$
(below $A$) has speed $v = \sqrt{2g|y|}$ where $g$ is gravitational
acceleration. The arc length element is $ds = \sqrt{1 + (y')^2} \,
dx$. The travel time element is $dt = ds/v$. The total time is:
\[
T[y] = \int_0^a \frac{\sqrt{1 + (y')^2}}{\sqrt{2g|y|}} \, dx.
\]
The problem is to find $y(x)$ minimizing $T[y]$ subject to $y(0) =
0$ and $y(a) = -b$.

\emph{Variational formulation.} The functional $T[y]$ is the
integral of a Lagrangian:
\[
T[y] = \int_0^a L(y, y') \, dx, \quad L(y, y') = \frac{\sqrt{1 +
(y')^2}}{\sqrt{2g|y|}}.
\]
The Lagrangian does not depend explicitly on $x$ (it depends only on
$y$ and $y'$); the corresponding Beltrami identity will simplify
the analysis.

\emph{Derivation of the Euler-Lagrange equation.} To find a curve
$y$ that extremizes $T[y]$, consider a perturbation $y + \epsilon
\eta$ where $\eta$ is a smooth function vanishing at the endpoints
($\eta(0) = \eta(a) = 0$) and $\epsilon$ is small. The first-order
variation of $T$ is:
\[
\frac{d}{d\epsilon}\bigg|_{\epsilon=0} T[y + \epsilon \eta] =
\int_0^a \left(\frac{\partial L}{\partial y} \eta + \frac{\partial
L}{\partial y'} \eta'\right) dx.
\]
Integration by parts on the second term, using $\eta(0) = \eta(a) =
0$:
\[
\int_0^a \frac{\partial L}{\partial y'} \eta' \, dx = -\int_0^a
\frac{d}{dx}\frac{\partial L}{\partial y'} \eta \, dx.
\]
Combining:
\[
\frac{d}{d\epsilon}\bigg|_{\epsilon=0} T = \int_0^a \left(\frac{\partial
L}{\partial y} - \frac{d}{dx}\frac{\partial L}{\partial y'}\right)
\eta \, dx.
\]
For $T$ to be extremized at $y$, this variation must vanish for all
admissible $\eta$. By the fundamental lemma of the calculus of
variations:
\[
\frac{\partial L}{\partial y} - \frac{d}{dx}\frac{\partial
L}{\partial y'} = 0.
\]
This is the Euler-Lagrange equation. It is a necessary condition for
$y$ to be an extremum of $T[y]$. The equation is a second-order
ordinary differential equation for $y(x)$.

\emph{Application to the brachistochrone.} For the brachistochrone
Lagrangian, the Beltrami identity (a first integral when $L$ is
independent of $x$) gives:
\[
L - y' \frac{\partial L}{\partial y'} = C
\]
for some constant $C$. Computing:
\[
\frac{\sqrt{1 + (y')^2}}{\sqrt{-2gy}} - \frac{y' \cdot y'}{\sqrt{-2gy}
\sqrt{1 + (y')^2}} = \frac{1}{\sqrt{-2gy}\sqrt{1 + (y')^2}} = C.
\]
Squaring: $-2gy(1 + (y')^2) = 1/C^2$, or $y(1 + (y')^2) = -1/(2gC^2)
= -k$ for some positive constant $k$.

\emph{Solving the equation.} The first integral is $y(1 + (y')^2) =
-k$. Substitute $y = -k \sin^2(\theta/2)$ (this is a clever
substitution motivated by the cycloid form). Then $y' = -k
\sin(\theta/2)\cos(\theta/2) \cdot \theta'/2 \cdot$ (something);
after careful computation, one finds:
\[
x = \frac{k}{2}(\theta - \sin\theta), \quad y = -\frac{k}{2}(1 -
\cos\theta).
\]
With $r = k/2$, this is the cycloid parameterization. The constant
$k$ (equivalently $r$) is fixed by requiring the curve to pass
through $B = (a, -b)$.

\emph{Verification.} The cycloid satisfies the Euler-Lagrange
equation. It passes through the endpoints. It is the unique smooth
curve doing so for given endpoints. The travel time can be computed
in closed form: $T = \pi \sqrt{r/g}$ for the cycloid arc from cusp
to first descent.

The brachistochrone is the chapter's central example. The result is
not merely that the cycloid is admissible; the result is that the
cycloid is forced. Among all admissible curves connecting $A$ and
$B$, the cycloid is the unique one that minimizes the travel time.
The selection is unique.

The Euler-Lagrange equation is the chapter's central tool. Given any
functional $J[y] = \int L(y, y', x) \, dx$, the equation
\[
\frac{\partial L}{\partial y} - \frac{d}{dx}\frac{\partial L}{\partial
y'} = 0
\]
is the necessary condition for $y$ to be an extremum of $J$. The
equation generalizes to multiple dependent variables (vector $y$,
producing a system of equations), to higher-order Lagrangians
(involving $y''$, producing fourth-order equations), and to multiple
independent variables (producing partial differential equations).

For mechanics, the Lagrangian is $L = T - V$ (kinetic minus potential
energy). The Euler-Lagrange equation gives Newton's equation $F =
ma$. The variational formulation and Newton's formulation are
equivalent; either implies the other. The variational form is more
fundamental: it generalizes to systems where Newton's form is
awkward (rotating frames, constrained motion, field theories).

For optics, Fermat's principle of least time gives the optical
equivalent: light follows the path of minimum travel time. The
Euler-Lagrange equation produces Snell's law and the lens equation.

For relativistic mechanics, the action $S = -mc^2 \int d\tau$ (mass
times the proper-time interval) is the Lagrangian; the
Euler-Lagrange equation produces the geodesic equation: free
particles follow geodesics of spacetime.

In Lean, the Euler-Lagrange equation appears in the
\texttt{GalerkinProcess} typeclass: the Galerkin projection of
Chapter 5 is the minimum-energy approximation; the
Euler-Lagrange equation is the variational form of the same
selection. The compiler's elaborator uses similar variational
ideas in its search heuristics: among admissible proofs, the
elaborator selects the one with minimum elaboration cost.

The chapter's claim is that the Euler-Lagrange equation is the
necessary condition for selection by a functional. Given the
admissible family and the functional, the equation produces the
selected member. The brachistochrone's cycloid is the canonical
example: a specific functional (the descent time) produces a
specific curve (the cycloid) through the Euler-Lagrange machinery.

\section{Kolmogorov Selection}
\label{se:kolmogorov-selection}

The Euler-Lagrange equation selects the extremum of a continuous
functional defined on an admissible family of curves. The selection
principle is geometric: among smooth curves through the boundary
data, the extremum minimizes or maximizes the functional. The
selection has a different character when the admissible family is
discrete or when the functional is the description length of the
candidate.

Kolmogorov complexity $K(w)$ of a string $w$ is the length of the
shortest program that produces $w$ as output. The complexity is
measured relative to a fixed universal Turing machine; the choice of
machine affects $K$ by only an additive constant, so the asymptotic
behavior is machine-independent.

Kolmogorov complexity is the canonical algorithmic measure of
informational content. A string of one million identical characters
(``aaaa...a'') has small $K$: a short program (``print `a' one
million times'') produces it. A string of one million random
characters has large $K$: no program much shorter than the string
itself produces it.

The chapter's third claim is that Kolmogorov complexity provides a
selection principle. Among admissible continuations of a given
record, the one with minimum $K$ is the most economical: it
asserts the least unnecessary structure. The selection is
analogous to the Euler-Lagrange selection (minimum of a
functional), but in the discrete algorithmic setting.

For the cubic spline of Chapter 5, the connection is direct. The
spline minimizes the integrated squared curvature $\int (f'')^2 \,
dx$. The minimum-curvature continuation of given data is also the
minimum-Kolmogorov-complexity continuation in the following sense:
among admissible $C^2$ continuations of the data, the spline has
the shortest description (the data plus the spline construction
algorithm) compared to alternatives that introduce extra
inflections.

This is informal because exact $K$ is not computable; it cannot be
verified by an algorithm that the spline is exactly the minimum-$K$
continuation. But the integrated curvature is a computable upper
bound on $K$: a function with more curvature has more inflection
points to describe, hence more bits in any reasonable encoding.

\begin{theorem}[Kolmogorov-Spline Correspondence (informal)]
For data $\{(x_i, y_i)\}_{i=0}^n$ and the admissible class of $C^2$
interpolants, the natural cubic spline approximately minimizes the
Kolmogorov complexity in the sense that any $C^2$ interpolant with
strictly less integrated curvature has strictly less complexity (up
to a constant accounting for the description algorithm).
\end{theorem}

This is not a precise theorem; it is a heuristic statement of the
informational economy that the spline embodies. The exact form
involves the choice of universal Turing machine and a careful
specification of what an interpolant's description includes.

The chapter takes Kolmogorov selection as a conceptual frame, not as
a computable algorithm. The frame says: among admissible
continuations, the one with the minimum informational content
(measured by some computable proxy for $K$) is the selected
continuation. The proxy varies with the situation: integrated
curvature for splines, action for mechanical paths, surface area
for soap films, edit distance for sequence alignments.

A critical observation: exact Kolmogorov complexity is not
computable. By Chaitin's theorem (a consequence of the Halting
Problem), no algorithm computes $K(w)$ for arbitrary $w$. The
selection principle ``minimize $K$'' is therefore not directly
implementable; it must be approximated.

The compiler's elaborator, when it selects a proof, uses a
computable proxy: the number of heartbeats (internal elaboration
steps) consumed by the proof. The heartbeat count is a computable
upper bound on the proof's algorithmic complexity; the elaborator
selects the proof with the lowest heartbeat count among admissible
proofs that succeed. This is a Kolmogorov-style selection with a
specific proxy.

\begin{leanexcerpt}{Proxy Cost}
structure EKG where
  private reference : Reading
  embigger? : Prop $\to$ Prop $\to$ Prop
\end{leanexcerpt}

The \texttt{EKG} structure in \texttt{LeanCalibration.lean}
encapsulates the heartbeat proxy. The \texttt{reference} field
carries the baseline reading (privately, hidden from user code).
The \texttt{embigger?} method is a binary predicate on
\texttt{Prop} that the calibration uses to certify one
heartbeat-bearing claim as no smaller than another. The proxy is
computable; the predicate is admissible; the selection of the
minimum is decidable through the calibration interface.

The chapter's discipline: when the exact Kolmogorov selection is
not implementable, use a computable proxy. The proxy should be a
plausible upper bound on $K$. The selection by the proxy is the
admissible substitute for the ideal selection.

The Kolmogorov frame extends to several specific selection problems
beyond the spline:

\emph{Sparse regression.} Given data and a large set of candidate
features, select the smallest set of features that explains the
data adequately. The Lasso regression solves this approximately by
minimizing the residual plus a penalty proportional to the
$\ell_1$ norm of the coefficients. The penalty is a computable
proxy for $K$.

\emph{Model selection.} Given several candidate statistical models,
choose the one that best balances goodness of fit against
complexity. The Bayesian information criterion (BIC) is a proxy
for $K$: BIC = $k \ln n - 2 \ln L$ where $k$ is the number of
parameters and $L$ is the likelihood. Selecting the model with
the smallest BIC is a computable approximation to Kolmogorov
selection.

\emph{Compression.} Given a string $w$, find the shortest encoding
that recovers $w$. Lempel-Ziv compression is the practical
algorithm; the resulting compressed length is a computable upper
bound on $K(w)$. The selection (among possible encodings) of the
shortest is the Kolmogorov principle applied to encoding.

Each of these is a Kolmogorov-style selection problem with a
specific computable proxy. The chapter's general principle holds:
among admissible alternatives, select the one with the minimum
informational content as measured by the chosen proxy.

In Lean, the \texttt{ChaitinsNumberSequence} in
\texttt{Episode3.lean} represents the halting probability $\Omega$,
the formal object whose digits encode the answers to the Halting
Problem. The compiler does not compute these digits; the type
exists but no values are constructed. The Kolmogorov selection at
the universe's level would require computing $\Omega$; since this
is impossible, the compiler uses heartbeats as a proxy.

The chapter's distinction: the Kolmogorov frame is a conceptual
ideal; the proxies are practical implementations. The
brachistochrone (a continuous selection problem) is solved exactly
by the Euler-Lagrange equation. The spline (a discrete selection)
is solved exactly by linear algebra. The general Kolmogorov
selection problem is unsolvable exactly; computable proxies are
the discipline.

The chapter's claim is that variational calculus (continuous,
Euler-Lagrange) and Kolmogorov selection (discrete, proxy-based)
are two faces of the same selection principle. Both ask: among
admissible alternatives, which has the minimum value of the
functional / complexity / energy? The answer is unique when the
functional is convex; for non-convex functionals or for
incomputable complexities, the selection is approximate, and the
proxies must be honest about their limitations.

\section{Fluxions Resolved}
\label{se:fluxions-resolved}

Newton's calculus was based on \emph{fluxions}: infinitesimal
quantities representing momentary changes. The derivative of $x$
with respect to $t$, written $\dot{x}$ in Newton's notation, was the
fluxion: the ratio of the infinitesimal change in $x$ to the
infinitesimal change in $t$ at a single instant.

The fluxion was philosophically problematic. Bishop Berkeley
famously criticized fluxions in \emph{The Analyst} (1734) as
``ghosts of departed quantities'': simultaneously zero (at the
limiting instant) and nonzero (a meaningful ratio). The
contradiction was not resolved until Cauchy's epsilon-delta
formalism in the early nineteenth century made the limit rigorous.

In the chapter's framework, the fluxion is resolved differently. The
derivative is not a ratio of infinitesimals; it is a Galerkin
limit of finite ratios. The discrete data are the counts; the
ratios are admissible because they are computed from finite
counts; the derivative is the limit of the ratios as the data are
refined.

For a function $x(t)$ recorded at discrete times $t_0 < t_1 < t_2 <
\ldots$, the finite difference $\Delta x_n / \Delta t_n =
(x(t_{n+1}) - x(t_n))/(t_{n+1} - t_n)$ is a ratio of two recorded
counts: the change in $x$ over the change in $t$. The finite
difference is admissible: both numerator and denominator are
counts; the ratio is well-defined.

As the time grid is refined ($\Delta t_n \to 0$), the finite
differences converge (if $x$ is smooth) to a limit. The limit is
the derivative $\dot{x}(t)$ in the modern sense. The limit is the
Galerkin projection of the ratio onto the increasingly fine grid:
the unique admissible continuation of the discrete ratios.

The chapter's resolution: fluxions are not infinitesimals. They are
limits of finite ratios. The infinitesimal vocabulary is shorthand
for the limit; the limit is admissible because it is the Galerkin
projection's limit value. No ghosts; no contradictions.

In Lean, derivatives are defined as limits via the type
\texttt{HasDerivAt}. The definition is the standard
epsilon-delta limit: $f$ has derivative $f'(a)$ at $a$ when for
every $\varepsilon > 0$, there exists $\delta > 0$ such that
$|f(x) - f(a) - f'(a)(x - a)| < \varepsilon |x - a|$ for all $|x
- a| < \delta$. The definition is constructive: the derivative is
computed from finite data (the values of $f$ near $a$) by taking
the limit.

The chapter's fourth claim is that calculus is grounded in finite
counting, not in infinitesimals. The derivative is the limit of
finite-difference ratios; the integral is the limit of
finite-sum Riemann approximations; the curvature is the limit of
finite second-difference ratios. All of calculus is admissible as
limits of counted operations.

This perspective unifies the chapter's themes. The Euler-Lagrange
equation involves $d/dx$ of $\partial L/\partial y'$. The
derivative $d/dx$ is the limit of finite differences; the
equation is admissible as a statement about Galerkin limits. The
selection by minimum-functional is an algebraic operation on the
discrete data; the limit (as the discretization refines) is the
continuous extremum.

The chapter's discipline: no infinitesimals; only limits of
finite operations. The notation $dx$ is shorthand; the meaning is
the limit. The selection principle is unaffected by the
notational simplification; the Euler-Lagrange equation is the
limit form of the discrete extremality condition.

\section{Galerkin Convergence}
\label{se:galerkin-convergence}

Consider the boundary-value problem $u^{(4)}(x) = f(x)$ on $[0, 1]$
with boundary conditions $u(0) = u(1) = u'(0) = u'(1) = 0$. This is
the equation for the bending of an elastic beam under a distributed
load $f$. The chapter's framework solves it by Galerkin projection.

Let $V_h$ be the space of cubic spline functions on $[0, 1]$ with
knot spacing $h$ that vanish at the endpoints and have zero
derivatives at the endpoints. As $h$ decreases (more knots), $V_h$
grows. The Galerkin projection $u_h$ is the unique solution in
$V_h$ of the variational form: $\int_0^1 u_h''(x) v''(x) \, dx =
\int_0^1 f(x) v(x) \, dx$ for all $v \in V_h$.

The Galerkin Convergence Theorem states that the sequence $u_h$
converges to the exact solution $u$ as $h \to 0$. The convergence
rate depends on the smoothness of $u$ and the order of the basis:
for cubic splines and $u \in H^4$, the rate is $\|u - u_h\|_{L^2}
\le C h^4 \|u^{(4)}\|_{L^2}$.

The chapter's fifth claim is that the Galerkin projection
converges to the exact solution as the discretization refines.
The discrete approximations are admissible at each finite
resolution; the limit is the continuous solution; the limit is
admissible because the Cauchy completion of the discrete
approximations is the continuum's element.

This is the chapter's continuous extension of Chapter 5's spline
result. Chapter 5 proved that the spline is the
minimum-curvature interpolant of finite data. The current
chapter extends this to differential equations: the spline
discretization of a variational problem converges to the
continuous solution as the discretization refines.

The convergence is the chapter's structural unification of the
discrete and the continuous. The discrete Galerkin solutions
$u_h$ are admissible at each finite resolution. The continuous
solution $u$ is the limit of the discrete approximations. The
limit is the continuous extremum of the same variational
problem. The discrete and continuous are unified through the
Galerkin convergence.

In Lean, the convergence theorem is part of \texttt{Mathlib}'s
functional-analysis library (for the underlying functional-analytic
results) plus the chapter's spline-specific machinery. The
\texttt{GalerkinProcess} typeclass requires the convergence proof
as part of its admissibility certificate. The compiler enforces
this: an instance must demonstrate convergence; without the
proof, the instance is rejected.

The chapter closes here. The selection principle (Section 1) chooses
the extremum of a functional. The Euler-Lagrange equation (Section
2) gives the necessary condition for the extremum. Kolmogorov
selection (Section 3) extends this to discrete settings via
computable proxies. Fluxions (Section 4) are resolved as limits of
finite ratios. Galerkin convergence (Section 5) unifies the
discrete and continuous through the convergence of finite
approximations to the continuous limit.

\begin{bridge}

The chapter's question was what selects one admissible continuation
as the history actually carried forward.

The answer is: minimization of a functional. The soap film selects
the minimum-area surface; the brachistochrone selects the
minimum-time path (the cycloid); the spline selects the
minimum-curvature interpolant; the Galerkin projection selects the
minimum-energy approximation. Each selection is forced by its
functional. The Euler-Lagrange equation is the necessary
condition: at the extremum, the functional derivative vanishes.
Kolmogorov selection extends the principle to discrete settings
with computable proxies for the incomputable exact $K$.

The chapter's discipline: among admissible continuations, the
minimum of the chosen functional is selected. The selection is
unique when the functional is strictly convex. Fluxions are
resolved as limits of finite ratios. Galerkin convergence unifies
the discrete and continuous through limit processes.

What remains is transport. The chapter has selected one
continuation; how does information move between continuations?
Chapter 7 takes up the transport question: how can structure
move from one record to another while preserving admissibility?

\end{bridge}

\begin{coda}{The Calligrapher Drawing One Stroke}

The calligrapher sits at a low wooden table in a quiet room. The
table is covered with a felt mat; on the mat is a sheet of
hand-made paper, smooth, slightly cream-colored, twelve centimeters
square. The brush rests in a small ceramic stand to the right.
The ink stone is to the left, with a small puddle of ink at its
center. The water dish is behind. The room is silent except for the
soft sound of the calligrapher's breath.

She has been practicing this stroke for thirty years. The stroke is
the horizontal first stroke of the character meaning ``one'': a
single line drawn from left to right, starting with a soft entry,
swelling slightly in the middle, and ending with a controlled tail.
The character is the simplest; the stroke is the most demanding.
The simplest character requires the most precise stroke.

She loads the brush. The brush is held vertically; the tip is dipped
into the ink puddle; the brush rotates slowly; the ink rises into
the hairs of the brush by capillary action. The calligrapher
withdraws the brush, taps it once gently on the side of the stone
to remove excess ink, and considers the load. The brush carries the
right amount: not too little (the stroke will skip), not too much
(the stroke will blot). The load is calibrated.

She positions the brush above the paper. The starting point of the
stroke is on the left side, about three centimeters from the
upper edge. The ending point is on the right side, three
centimeters from the upper edge, parallel to the starting point.
Between these two points, the stroke must travel.

She pauses. The pause is not idle. The calligrapher is selecting
the stroke. Not consciously; her training has internalized the
selection. But the selection is happening: among all the possible
strokes connecting the two points, which is the one she will draw?

The answer is the stroke of minimum added curvature: the stroke
that introduces no inflection not forced by the boundary data and
the brush's natural behavior. The selection is the chapter's
selection principle, performed by the calligrapher's body. The
brush is the instrument; the paper is the substrate; the
minimization is the discipline.

The minimum-curvature constraint is the calligrapher's spline. The
brush will travel from start to end; the brush's bristles will
flex slightly as the stroke proceeds; the paper will absorb the
ink slightly; all of these are physical processes that the
calligrapher must accommodate. The optimal stroke is the
trajectory of the brush that minimizes unnecessary effort while
producing the intended visual result.

She begins. The brush descends. The first contact is the entry:
the tip of the brush touches the paper, the bristles spread
slightly, the first dot of ink appears. The entry is a controlled
event: the calligrapher's wrist initiates the contact at a
specific angle, with a specific pressure, with a specific speed.
The control determines the entry's appearance: a soft entry, not
a hard impact.

The brush moves to the right. The speed is constant; the
pressure varies; the ink flow varies; the bristles flex and
recover. The calligrapher's arm carries the brush horizontally;
her wrist makes small adjustments to maintain the stroke's
character. The stroke is being drawn.

This is the Euler-Lagrange equation in physical form. The
functional is the calligrapher's aesthetic: the visual quality of
the stroke. The admissible curves are the trajectories the brush
could take from start to end. The selection is the trajectory
that the calligrapher's training has identified as the unique
admissible answer. The training is the calibration; the stroke
is the executed continuation.

The middle of the stroke is the swell: a slight thickening as the
brush carries its full ink and the bristles flex maximally. The
swell is not a deliberate addition; it is the natural consequence
of the brush's full engagement with the paper. The calligrapher
does not add the swell; she allows the brush's natural behavior
to produce it. The minimization principle accommodates this:
the swell is what minimum-effort produces, not what
maximum-effort produces.

The stroke continues. The brush approaches the end. The
calligrapher decreases the pressure; the bristles return toward
their resting shape; the stroke tapers. The taper is the exit:
a controlled lift of the brush from the paper, producing a
sharp tail. The taper is the mirror of the entry: same precision,
same controlled mechanics, opposite direction.

The brush lifts. The stroke is complete. The horizontal line is
on the paper: an inked trajectory from left to right, with a
soft entry, a natural swell, and a controlled tail. The total
duration was perhaps three seconds.

The calligrapher inspects the stroke. The result is acceptable.
The line is straight; the swell is in the right place; the
tail is correct. The selection has succeeded: the minimum-effort
stroke has produced the intended visual quality.

What if the stroke had failed? Suppose the brush had wavered, or
the ink had blotted, or the line had been slightly curved. The
failure would have been an admission that the selection had not
quite produced the unique minimum-effort answer. The
calligrapher's training would identify the source of the failure
(too much ink, too quick a motion, too high a brush angle) and
correct for the next stroke.

The chapter's discipline is what the calligrapher embodies. The
admissible continuations are the possible stroke trajectories.
The functional is the aesthetic / mechanical / inking
requirements. The Euler-Lagrange equation is the implicit
calculation that the calligrapher's body performs to identify
the minimum-functional trajectory. The execution is the brush's
actual motion. The result is the inked line on the paper.

The line on the paper is the chapter's structures performed in
ink. The boundary data are the start and end points. The
admissible continuations are the possible strokes. The selection
is the minimum-curvature stroke. The execution is the brush's
trajectory. The residue is the slight imperfection that any
physical execution introduces.

The calligrapher will draw another stroke now. The next character
in her practice sheet is the character meaning ``two'': two
horizontal strokes, the lower one slightly longer. Each stroke
is its own selection problem; each admits its own
Euler-Lagrange solution; each is executed by the same
discipline. After thirty years, the calligrapher's bodily
training is the unconscious implementation of the chapter's
variational principle.

She loads the brush again. The next stroke begins.

\end{coda}
